% !TeX root = ../papers/nc-p3.tex
\section{多項式成長による一般倍率間の第一拡張と二段積の制限}
\label{sec:polynomial-growth-extension-bridges}
以下では折り畳み $B$ の二次表示と多項式成長を直接使う。
拡張は右加群の上の層から下の層へ書く。$S_u$ は非頂点単純 $S$ の矢の次数倍率 $u$ による移動で、
非零な正長語の跡により同型となる倍率は有限個である。

\begin{theorem}\label{thm:growth-generic-extensions}
標準正則窓と有限次元非頂点単純 $S$ に対し、一般の $u,v$ で
$\operatorname{Ext}^1_B(S_u,S_v)=0$ である。また頂点単純を経由する積
\[
 \bigoplus_j\operatorname{Ext}^1(S_u,k_j)\otimes\operatorname{Ext}^1(k_j,S_v)
 \longrightarrow\operatorname{Ext}^2(S_u,S_v)
\]
は一般の $u,v$ で単射である。従って $d=\dim S$ とすれば
\[
 \sum_j a_j b_j\le\chi(d,d)
 =4(d_0+d_2-d_1-d_3)^2-2\big((d_0-d_2)^2+(d_1-d_3)^2\big).
\]
有限辺も行列恒等式も仮定しない。
\end{theorem}

\begin{proof}[第一拡張と成長の比較]
一般に第一Extが正なら、任意の $n$ に対し相異なる $S_1,\ldots,S_n$ を取り、
全順序対 $i\ne j$ に非内部な相対導分 $C_{ij}$ を取れる。
実際Extの次元は行列の階数で決まり、各対に要求する開条件の有限交叉を取ればよい。
頂点冪等元に相対的な矢行列で記述する。
頂点 $1,\ldots,n$ の完全有向グラフの経路代数を係数に、
各矢を対角の元の表現と非対角 $C_{ij}x_{ij}$ の和に送る。
ここで $x_{ij}$ は $i$ から $j$ への文字であり、合成は経路順とする。
元の六関係の各位置での定数と一次係数は既に零である。残る二次行列の各成分を
関係として係数環を割る。こうして $B$ の表現ができる。係数環 $D$ は経路長で次数付きで、
順序対あたりの関係数は $n$ に依らない定数 $K$ 以下である。

ここで使う成長判定を記す。有限箙の矢数行列を $A$、次数 $s\ge2$ の関係数行列を $R_s$、
$H_{ij}(t)=\sum_m\dim(e_iD_me_j)t^m$ とする。各終点への経路を最後の矢で分けると
\[
 H(t)\big(I-tA+\textstyle\sum_s t^s R_s\big)\ \ge\ I
\]
が係数ごとに成り立つ。最後の矢を外す表示の核は、末尾に関係を置いたベクトルで左生成されるためである。
表示する関係の独立性は不要である。ある $0<t<1$ と全成分正の $h$ で括弧が
$h$ を成分ごとに負に送れば $H(t)$ は収束できない。この場合 $D$ は指数成長する。
今は $h_i=1$, $t=2/n$ で行和は $1-2(n-1)/n+4K/n$ 以下となる。

この成長は $B$ の有限個の元の積でも検出できる。次数2以上を切った表現の像を取ると、
定数商は $\prod_i\operatorname{End}(S_i)$ である。一次の核は二乗零の両側加群で、
各 $i,j$ 矩形の重複度は1である。ある矩形への射影が零ならそこに制限した拡張は
半単純代数を通り分裂し、$C_{ij}$ の選択に反する。従って核は全ての非対角矩形を含む。
各頂点で選んだ一つの座標間の行列単位に $x_{ij}$ を掛けたものを、
一次の先頭項に持つ $B$ の元を取る。それらの積の先頭次数は対応する経路を成分に持つ
（経路が関係で零ならその次数の成分は零）。各次数で経路の基底を選べるので、
有限個の持上げの最大語長を $L$ として $B$ の長さ $Lm$ 以下の空間は
$\dim D_m$ 以上である。多項式成長に反する。
\end{proof}

\begin{lemma}\label{lem:growth-fixed-middle}
$Q_1,\ldots,Q_r$ を組成因子が全て頂点単純で次数倍率で同型類が動かない固定有限次元加群とする。
一般の $u,v$ に対し、積で零に行くテンソル
\[
 \tau\in\bigoplus_l \operatorname{Ext}^1(S_u,Q_l)\otimes\operatorname{Ext}^1(Q_l,S_v)
\]
を階数分解し、中間に各 $Q_l$ の対応する個数の直和を入れた三段加群を作る。
その有限像において $S_u$ のブロックから頂点型のブロックを経て $S_v$ のブロックへ至る
非零積を持つような $\tau$ は一般の対には存在しない。
ここでブロックは半単純商の冪等元の直交持上げで切り出す。
\end{lemma}
\begin{proof}
そのような対が稠密にあれば一般の $n$ 個の相異なる倍率で全順序対 $i\ne j$ が条件を満たす。
実際拡張、階数分解、三段への持上げ、像のブロックと非零積の条件は全て構成可能である。
有限像を張る語長はその周囲の行列空間の次元で抑えられる。
$a_l=\dim\operatorname{Ext}^1(S_i,Q_l)$、$b_l=\dim\operatorname{Ext}^1(Q_l,S_j)$
と置く。倍率での同型性により一定で、$a_l b_l=0$ の添字は落とす。
$n$ 個の端点と $r$ 個の中間点を持つ箙に、各 $i$ から $l$ へ $a_l$ 本の矢を、
各 $l$ から $j$ へ $b_l$ 本の矢を置く。端点間の長さ2の経路は各対で $P=\sum_l a_l b_l$ 本ある。
$i\ne j$ ではそれらを $\tau_{ij}$ の各係数倍の一文字 $z_{ij}$ に同一視する
（非零係数の経路の定数倍を $z_{ij}$ とする）。対角では全て零にする。
これは非対角 $P-1$、対角 $P$ 個の二次関係である。

対角に $S_i,Q_l$ の矢行列、一次に各Ext基底の導分を置き、さらに $i,j$ 成分に
二次補正 $F_{ij}z_{ij}$ を置く。元の関係の $i,j$ 成分の二次係数は補正前には
$\tau_{ij}$ のカップ積を表し、零類なので全関係で同時にこれを打ち消す $F_{ij}$ を選べる。
これは二次関係に上三角三段行列を代入した最後の非対角式そのものである。
中間点間に残る二次の元の関係、全成分の三次、四次係数をそれぞれ斉次関係として加える。
これで得た係数環 $D$ 上に $B$ が作用する。各追加次数の関係数は各順序対あたり
一定の $K$ 以下である。端点間の二次に追加はなく、三次は異種の端点間、四次は両端が $S$ 型の場合だけである。

$t=c/\sqrt n$, $h_i=1$, $h_l=\sqrt n(c b_l-\epsilon)$ と取る。
$c>1$ を固定し、全中間重みが正で $1-c^2+c\epsilon\sum_l a_l<0$ となるよう
$\epsilon>0$ を固定する。上の成長判定の端点行の上界は
\[
 1-c^2P+c\epsilon\sum_l a_l+c^2(P-1)+O(n^{-1}),
\]
中間行は $-\epsilon\sqrt n+O(n^{-1/2})$ である。例えば端点から中間への三次関係は
$rKt^3\max_l h_l=O(n^{-1})$、中間から端点への三次は $nKt^3=O(n^{-1/2})$ である。
従って大きい $n$ で $D$ は指数成長する。

先頭次数を移す際、中間加群自身は半単純とは限らない。次数3以上を切った表現の像 $C$ は
定数像へ冪零核で全射する。この定数像は $S_i$ の各全行列環と全中間での像との直積である
（その単純商の型が交わらず、各 $S_i$ は一回だけ現れる）。
$S_i$ の行列単位と残りの単位 $e_Q$ を定数像から $C$ に持ち上げる。
一次まででは $e_i C e_Q$ は $i$ から中間への矩形だけに台を持ち、
$e_Q C e_j$ も中間から $j$ への矩形だけに台を持つ。
実際これらは定数項が零で、左右の持上げの定数部でそれぞれの先頭項を切れる。
三段表現の $i,j$ 成分の二次係数を見るには、その対の経路だけを三角形の三層に作用させればよい。
中間層の重複空間への一次矢は $\tau_{ij}$ の階数分解の行と列で作用し、
長さ二の合成はちょうど $z_{ij}$ の係数となる。他の二次以上の関係はこの三層で自動的に零になる。
従ってこの三段像は $C$ の商であり、非零積の仮定から
$e_i C e_Q C e_j$ の像は非零である。持上げの選び方は根基による同時共役しか変えず、
$F_{ij}$ の選び方も二つの一次矩形の積を変えない。
左右を行列単位で切れば、指定座標間の $z_{ij}$ だけを二次先頭項に持つ $B$ の元を得る。
両端が端点の全偶数長経路はこれら $z_{ij}$ の積で張られ、その他の経路も両端の矢を
高々一本ずつ外せばこの形になる。有限個の矢しかないのでこの部分も指数成長する。
第一拡張の場合と同じ語長比較が矛盾を与える。
\end{proof}

\begin{proof}[定理の二段積とEuler式]
中間が互いに異なる頂点単純なら、非零テンソルの階数分解は各非零成分で両側に独立な拡張を与える。
補題\ref{lem:pi-yoneda-bridges}の矩形の証明により、積が零なら三段像においてその頂点型を
通る非零積がある。補題\ref{lem:growth-fixed-middle}がこれを排除する。
相異なる一般倍率ではHomと双対の第四Ext、両向きの第一Extと双対の第三Extが零。
折り畳みのKoszul複体のEuler式を展開すれば主張の数式を得る。
\end{proof}
